
\documentclass{article}
\usepackage{colortbl}
\usepackage{makecell}
\usepackage{multirow}
\usepackage{supertabular}

\begin{document}

\newcounter{utterance}

\centering \large Interaction Transcript for game `codenames', experiment `association\_difficult', episode 1 with gpt{-}5.4{-}azure.
\vspace{24pt}

{ \footnotesize  \setcounter{utterance}{1}
\setlength{\tabcolsep}{0pt}
\begin{supertabular}{c@{$\;$}|p{.15\linewidth}@{}p{.15\linewidth}p{.15\linewidth}p{.15\linewidth}p{.15\linewidth}p{.15\linewidth}}
    \# & $\;$A & \multicolumn{4}{c}{Game Master} & $\;\:$B\\
    \hline

    \theutterance \stepcounter{utterance}  
    & & & \multicolumn{2}{p{0.3\linewidth}}{
        \cellcolor[rgb]{0.9,0.9,0.9}{
            \makecell[{{p{\linewidth}}}]{
                \texttt{\tiny{[GM$|$GM]}}
                {'hidden': {'team': ['srebrny', 'irys', 'spódnica', 'chata', 'tenis', 'sprzeniewierzenie', 'jurta', 't-shirt', 'narcyz'], 'innocent': ['sufit', 'gruźlica', 'bez', 'jaskinia', 'pokój', 'jeansy', 'dżuma'], 'opponent': ['ściana', 'petunia', 'biały', 'przekroczenie prędkości', 'samochód', 'boks', 'dach', 'majtki'], 'assassin': ['porwanie']}, 'revealed': {'team': {'team': [], 'innocent': [], 'opponent': [], 'assassin': []}, 'opponent': {'team': [], 'innocent': [], 'opponent': [], 'assassin': []}}}
            }
        }
    }
    & & \\ \\

    \theutterance \stepcounter{utterance}  
    & & \multicolumn{4}{p{0.6\linewidth}}{
        \cellcolor[rgb]{0.9,0.9,0.9}{
            \makecell[{{p{\linewidth}}}]{
                \texttt{\tiny{[P1$\langle$GM]}}
                \texttt{Podaj jedną wskazówkę składającą się z pojedynczego słowa, która odnosi się do niektórych słów twojej drużyny z poniższej listy. Możesz wybrać, do ilu i do jakich słów odnosi się twoja wskazówka.} \\
\texttt{Celem jest, aby twój partner odgadł wszystkie słowa twojej drużyny na podstawie podanych wskazówek, zanim drużyna przeciwna odkryje wszystkie swoje. Jeśli twój partner odgadnie słowo 'przeciwnika' lub 'rozpraszające', jego tura zgadywania kończy się i drużyna przeciwna wykonuje turę. Jeśli twój partner odgadnie słowo 'zabójcze', gra kończy się i twoja drużyna natychmiast przegrywa.} \\
\texttt{Im więcej słów odnosi się do twojej wskazówki, tym szybciej możesz wygrać grę, ale tym trudniejsze mogą być skojarzenia do odgadnięcia dla twojego partnera, więc wybieraj słowo wskazówki mądrze. Słowo wskazówki musi być semantycznie powiązane ze słowami docelowymi, nie może być jednym ze słów na listach ani zawierać ich części.} \\
\\ 
\texttt{Zawsze podawaj swoją wskazówkę składającą się z pojedynczego słowa i oddzieloną przecinkami listę powiązanych słów docelowych w następującym formacie i staraj się, aby twoje odpowiedzi były jak najkrótsze, nigdy nie dołączaj żadnego innego tekstu niż jest to wymagane w tym formacie:} \\
\\ 
\texttt{WSKAZÓWKA: <SŁOWO>} \\
\texttt{CELE: <CELE>} \\
\\ 
\texttt{Słowa twojej drużyny to: srebrny, irys, spódnica, chata, tenis, sprzeniewierzenie, jurta, t{-}shirt, narcyz.} \\
\texttt{Słowa drużyny przeciwnika to: ściana, petunia, biały, przekroczenie prędkości, samochód, boks, dach, majtki.} \\
\texttt{Słowa rozpraszające to: sufit, gruźlica, bez, jaskinia, pokój, jeansy, dżuma.} \\
\texttt{Słowa zabójcze to: porwanie.} \\
            }
        }
    }
    & \\ \\

    \theutterance \stepcounter{utterance}  
    & \multicolumn{4}{p{0.6\linewidth}}{
        \cellcolor[rgb]{0.9,0.9,0.9}{
            \makecell[{{p{\linewidth}}}]{
                \texttt{\tiny{[P1$\rangle$GM]}}
                \texttt{WSKAZÓWKA: kort} \\
\texttt{CELE: tenis} \\
            }
        }
    }
    & & \\ \\

    \theutterance \stepcounter{utterance}  
    & & & \multicolumn{2}{p{0.3\linewidth}}{
        \cellcolor[rgb]{0.9,0.9,0.9}{
            \makecell[{{p{\linewidth}}}]{
                \texttt{\tiny{[GM$|$GM]}}
                \texttt{kort} \\
            }
        }
    }
    & & \\ \\

    \theutterance \stepcounter{utterance}  
    & & & \multicolumn{2}{p{0.3\linewidth}}{
        \cellcolor[rgb]{0.9,0.9,0.9}{
            \makecell[{{p{\linewidth}}}]{
                \texttt{\tiny{[GM$|$GM]}}
                ['tenis']
            }
        }
    }
    & & \\ \\

    \theutterance \stepcounter{utterance}  
    & & & \multicolumn{2}{p{0.3\linewidth}}{
        \cellcolor[rgb]{0.9,0.9,0.9}{
            \makecell[{{p{\linewidth}}}]{
                \texttt{\tiny{[GM$|$GM]}}
                {'word': 'tenis', 'assignment': 'team'}
            }
        }
    }
    & & \\ \\

    \theutterance \stepcounter{utterance}  
    & & \multicolumn{4}{p{0.6\linewidth}}{
        \cellcolor[rgb]{0.9,0.9,0.9}{
            \makecell[{{p{\linewidth}}}]{
                \texttt{\tiny{[GM$\rangle$P2]}}
                \texttt{Podaj oddzieloną przecinkami listę maksymalnie 1 słów z poniższej listy, które najlepiej odnoszą się lub są najbardziej związane ze słowem 'kort'. Zawsze rozpoczynaj swoją listę odgadnięć od ZGADNIĘCIE:  i nie dołączaj żadnego innego tekstu w swojej odpowiedzi.} \\
\\ 
\texttt{jaskinia, boks, majtki, dach, ściana, tenis, porwanie, pokój, chata, jeansy, jurta, sufit, biały, irys, sprzeniewierzenie, przekroczenie prędkości, gruźlica, narcyz, bez, srebrny, petunia, dżuma, samochód, t{-}shirt, spódnica} \\
            }
        }
    }
    & \\ \\

    \theutterance \stepcounter{utterance}  
    & & & \multicolumn{4}{p{0.6\linewidth}}{
        \cellcolor[rgb]{0.9,0.9,0.9}{
            \makecell[{{p{\linewidth}}}]{
                \texttt{\tiny{[GM$\langle$P2]}}
                \texttt{ZGADNIĘCIE: tenis} \\
            }
        }
    }
     \\ \\

    \theutterance \stepcounter{utterance}  
    & & & \multicolumn{2}{p{0.3\linewidth}}{
        \cellcolor[rgb]{0.9,0.9,0.9}{
            \makecell[{{p{\linewidth}}}]{
                \texttt{\tiny{[GM$|$GM]}}
                ['tenis']
            }
        }
    }
    & & \\ \\

    \theutterance \stepcounter{utterance}  
    & & & \multicolumn{2}{p{0.3\linewidth}}{
        \cellcolor[rgb]{0.9,0.9,0.9}{
            \makecell[{{p{\linewidth}}}]{
                \texttt{\tiny{[GM$|$GM]}}
                {'word': 'tenis', 'assignment': 'team'}
            }
        }
    }
    & & \\ \\

    \theutterance \stepcounter{utterance}  
    & & & \multicolumn{2}{p{0.3\linewidth}}{
        \cellcolor[rgb]{0.9,0.9,0.9}{
            \makecell[{{p{\linewidth}}}]{
                \texttt{\tiny{[GM$|$GM]}}
                {'word': 'tenis', 'assignment': 'team'}
            }
        }
    }
    & & \\ \\

    \theutterance \stepcounter{utterance}  
    & & & \multicolumn{2}{p{0.3\linewidth}}{
        \cellcolor[rgb]{0.9,0.9,0.9}{
            \makecell[{{p{\linewidth}}}]{
                \texttt{\tiny{[GM$|$GM]}}
                {'word': 'przekroczenie prędkości', 'assignment': 'opponent'}
            }
        }
    }
    & & \\ \\

    \theutterance \stepcounter{utterance}  
    & & & \multicolumn{2}{p{0.3\linewidth}}{
        \cellcolor[rgb]{0.9,0.9,0.9}{
            \makecell[{{p{\linewidth}}}]{
                \texttt{\tiny{[GM$|$GM]}}
                {'hidden': {'team': ['srebrny', 'irys', 'spódnica', 'chata', 'sprzeniewierzenie', 'jurta', 't-shirt', 'narcyz'], 'innocent': ['sufit', 'gruźlica', 'bez', 'jaskinia', 'pokój', 'jeansy', 'dżuma'], 'opponent': ['ściana', 'petunia', 'biały', 'samochód', 'boks', 'dach', 'majtki'], 'assassin': ['porwanie']}, 'revealed': {'team': {'team': ['tenis'], 'innocent': [], 'opponent': [], 'assassin': []}, 'opponent': {'team': [], 'innocent': [], 'opponent': ['przekroczenie prędkości'], 'assassin': []}}}
            }
        }
    }
    & & \\ \\

    \theutterance \stepcounter{utterance}  
    & & \multicolumn{4}{p{0.6\linewidth}}{
        \cellcolor[rgb]{0.9,0.9,0.9}{
            \makecell[{{p{\linewidth}}}]{
                \texttt{\tiny{[P1$\langle$GM]}}
                \texttt{Słowo tenis zostało odgadnięte prawidłowo. Tura twojego kolegi z drużyny się tam skończyła.} \\
\texttt{Teraz podaj kolejną wskazówkę odnoszącą się do niektórych z pozostałych słów twojego zespołu oraz listę powiązanych słów docelowych. Pamiętaj, aby rozpocząć wskazówkę od 'WSKAZÓWKA: ', wstawić nową linię i rozpocząć listę słów docelowych oddzielonych przecinkami od 'CELE: '. Uwaga: niektóre słowa zostały usunięte z list w porównaniu do poprzednich żądań.} \\
\\ 
\texttt{Twoje pozostałe słowa zespołu to: srebrny, irys, spódnica, chata, sprzeniewierzenie, jurta, t{-}shirt, narcyz.} \\
\texttt{Pozostałe słowa twojego przeciwnika to: ściana, petunia, biały, przekroczenie prędkości, samochód, boks, dach, majtki.} \\
\texttt{Pozostałe słowa rozpraszające to: sufit, gruźlica, bez, jaskinia, pokój, jeansy, dżuma.} \\
\texttt{Pozostałe słowa zabójcy to: porwanie.} \\
            }
        }
    }
    & \\ \\

    \theutterance \stepcounter{utterance}  
    & \multicolumn{4}{p{0.6\linewidth}}{
        \cellcolor[rgb]{0.9,0.9,0.9}{
            \makecell[{{p{\linewidth}}}]{
                \texttt{\tiny{[P1$\rangle$GM]}}
                \texttt{WSKAZÓWKA: defraudacja} \\
\texttt{CELE: sprzeniewierzenie} \\
            }
        }
    }
    & & \\ \\

    \theutterance \stepcounter{utterance}  
    & & & \multicolumn{2}{p{0.3\linewidth}}{
        \cellcolor[rgb]{0.9,0.9,0.9}{
            \makecell[{{p{\linewidth}}}]{
                \texttt{\tiny{[GM$|$GM]}}
                \texttt{defraudacja} \\
            }
        }
    }
    & & \\ \\

    \theutterance \stepcounter{utterance}  
    & & & \multicolumn{2}{p{0.3\linewidth}}{
        \cellcolor[rgb]{0.9,0.9,0.9}{
            \makecell[{{p{\linewidth}}}]{
                \texttt{\tiny{[GM$|$GM]}}
                ['sprzeniewierzenie']
            }
        }
    }
    & & \\ \\

    \theutterance \stepcounter{utterance}  
    & & & \multicolumn{2}{p{0.3\linewidth}}{
        \cellcolor[rgb]{0.9,0.9,0.9}{
            \makecell[{{p{\linewidth}}}]{
                \texttt{\tiny{[GM$|$GM]}}
                {'word': 'sprzeniewierzenie', 'assignment': 'team'}
            }
        }
    }
    & & \\ \\

    \theutterance \stepcounter{utterance}  
    & & \multicolumn{4}{p{0.6\linewidth}}{
        \cellcolor[rgb]{0.9,0.9,0.9}{
            \makecell[{{p{\linewidth}}}]{
                \texttt{\tiny{[GM$\rangle$P2]}}
                \texttt{Słowo tenis zostało odgadnięte prawidłowo. Twoja tura się tam skończyła.} \\
\texttt{Teraz podaj kolejną listę oddzieloną przecinkami, zawierającą co najmniej 1 i maksymalnie 1 słów z poniższej listy słów, które najlepiej odnoszą się lub są najbardziej związane ze słowem 'defraudacja'. Pamiętaj, aby rozpocząć odpowiedź od ZGADNIĘCIE: . Uwaga: niektóre słowa zostały usunięte z listy w porównaniu do poprzednich próśb.} \\
\\ 
\texttt{jaskinia, boks, majtki, dach, ściana, porwanie, pokój, chata, jeansy, jurta, sufit, biały, irys, sprzeniewierzenie, gruźlica, narcyz, bez, srebrny, petunia, dżuma, samochód, t{-}shirt, spódnica} \\
            }
        }
    }
    & \\ \\

    \theutterance \stepcounter{utterance}  
    & & & \multicolumn{4}{p{0.6\linewidth}}{
        \cellcolor[rgb]{0.9,0.9,0.9}{
            \makecell[{{p{\linewidth}}}]{
                \texttt{\tiny{[GM$\langle$P2]}}
                \texttt{ZGADNIĘCIE: sprzeniewierzenie} \\
            }
        }
    }
     \\ \\

    \theutterance \stepcounter{utterance}  
    & & & \multicolumn{2}{p{0.3\linewidth}}{
        \cellcolor[rgb]{0.9,0.9,0.9}{
            \makecell[{{p{\linewidth}}}]{
                \texttt{\tiny{[GM$|$GM]}}
                ['sprzeniewierzenie']
            }
        }
    }
    & & \\ \\

    \theutterance \stepcounter{utterance}  
    & & & \multicolumn{2}{p{0.3\linewidth}}{
        \cellcolor[rgb]{0.9,0.9,0.9}{
            \makecell[{{p{\linewidth}}}]{
                \texttt{\tiny{[GM$|$GM]}}
                {'word': 'sprzeniewierzenie', 'assignment': 'team'}
            }
        }
    }
    & & \\ \\

    \theutterance \stepcounter{utterance}  
    & & & \multicolumn{2}{p{0.3\linewidth}}{
        \cellcolor[rgb]{0.9,0.9,0.9}{
            \makecell[{{p{\linewidth}}}]{
                \texttt{\tiny{[GM$|$GM]}}
                {'word': 'sprzeniewierzenie', 'assignment': 'team'}
            }
        }
    }
    & & \\ \\

    \theutterance \stepcounter{utterance}  
    & & & \multicolumn{2}{p{0.3\linewidth}}{
        \cellcolor[rgb]{0.9,0.9,0.9}{
            \makecell[{{p{\linewidth}}}]{
                \texttt{\tiny{[GM$|$GM]}}
                {'word': 'petunia', 'assignment': 'opponent'}
            }
        }
    }
    & & \\ \\

    \theutterance \stepcounter{utterance}  
    & & & \multicolumn{2}{p{0.3\linewidth}}{
        \cellcolor[rgb]{0.9,0.9,0.9}{
            \makecell[{{p{\linewidth}}}]{
                \texttt{\tiny{[GM$|$GM]}}
                {'hidden': {'team': ['srebrny', 'irys', 'spódnica', 'chata', 'jurta', 't-shirt', 'narcyz'], 'innocent': ['sufit', 'gruźlica', 'bez', 'jaskinia', 'pokój', 'jeansy', 'dżuma'], 'opponent': ['ściana', 'biały', 'samochód', 'boks', 'dach', 'majtki'], 'assassin': ['porwanie']}, 'revealed': {'team': {'team': ['tenis', 'sprzeniewierzenie'], 'innocent': [], 'opponent': [], 'assassin': []}, 'opponent': {'team': [], 'innocent': [], 'opponent': ['przekroczenie prędkości', 'petunia'], 'assassin': []}}}
            }
        }
    }
    & & \\ \\

    \theutterance \stepcounter{utterance}  
    & & \multicolumn{4}{p{0.6\linewidth}}{
        \cellcolor[rgb]{0.9,0.9,0.9}{
            \makecell[{{p{\linewidth}}}]{
                \texttt{\tiny{[P1$\langle$GM]}}
                \texttt{Słowo sprzeniewierzenie zostało odgadnięte prawidłowo. Tura twojego kolegi z drużyny się tam skończyła.} \\
\texttt{Teraz podaj kolejną wskazówkę odnoszącą się do niektórych z pozostałych słów twojego zespołu oraz listę powiązanych słów docelowych. Pamiętaj, aby rozpocząć wskazówkę od 'WSKAZÓWKA: ', wstawić nową linię i rozpocząć listę słów docelowych oddzielonych przecinkami od 'CELE: '. Uwaga: niektóre słowa zostały usunięte z list w porównaniu do poprzednich żądań.} \\
\\ 
\texttt{Twoje pozostałe słowa zespołu to: srebrny, irys, spódnica, chata, jurta, t{-}shirt, narcyz.} \\
\texttt{Pozostałe słowa twojego przeciwnika to: ściana, petunia, biały, samochód, boks, dach, majtki.} \\
\texttt{Pozostałe słowa rozpraszające to: sufit, gruźlica, bez, jaskinia, pokój, jeansy, dżuma.} \\
\texttt{Pozostałe słowa zabójcy to: porwanie.} \\
            }
        }
    }
    & \\ \\

    \theutterance \stepcounter{utterance}  
    & \multicolumn{4}{p{0.6\linewidth}}{
        \cellcolor[rgb]{0.9,0.9,0.9}{
            \makecell[{{p{\linewidth}}}]{
                \texttt{\tiny{[P1$\rangle$GM]}}
                \texttt{WSKAZÓWKA: mongolska} \\
\texttt{CELE: jurta} \\
            }
        }
    }
    & & \\ \\

    \theutterance \stepcounter{utterance}  
    & & & \multicolumn{2}{p{0.3\linewidth}}{
        \cellcolor[rgb]{0.9,0.9,0.9}{
            \makecell[{{p{\linewidth}}}]{
                \texttt{\tiny{[GM$|$GM]}}
                \texttt{mongolska} \\
            }
        }
    }
    & & \\ \\

    \theutterance \stepcounter{utterance}  
    & & & \multicolumn{2}{p{0.3\linewidth}}{
        \cellcolor[rgb]{0.9,0.9,0.9}{
            \makecell[{{p{\linewidth}}}]{
                \texttt{\tiny{[GM$|$GM]}}
                ['jurta']
            }
        }
    }
    & & \\ \\

    \theutterance \stepcounter{utterance}  
    & & & \multicolumn{2}{p{0.3\linewidth}}{
        \cellcolor[rgb]{0.9,0.9,0.9}{
            \makecell[{{p{\linewidth}}}]{
                \texttt{\tiny{[GM$|$GM]}}
                {'word': 'jurta', 'assignment': 'team'}
            }
        }
    }
    & & \\ \\

    \theutterance \stepcounter{utterance}  
    & & \multicolumn{4}{p{0.6\linewidth}}{
        \cellcolor[rgb]{0.9,0.9,0.9}{
            \makecell[{{p{\linewidth}}}]{
                \texttt{\tiny{[GM$\rangle$P2]}}
                \texttt{Słowo sprzeniewierzenie zostało odgadnięte prawidłowo. Twoja tura się tam skończyła.} \\
\texttt{Teraz podaj kolejną listę oddzieloną przecinkami, zawierającą co najmniej 1 i maksymalnie 1 słów z poniższej listy słów, które najlepiej odnoszą się lub są najbardziej związane ze słowem 'mongolska'. Pamiętaj, aby rozpocząć odpowiedź od ZGADNIĘCIE: . Uwaga: niektóre słowa zostały usunięte z listy w porównaniu do poprzednich próśb.} \\
\\ 
\texttt{jaskinia, boks, majtki, dach, ściana, porwanie, pokój, chata, jeansy, jurta, sufit, biały, irys, gruźlica, narcyz, bez, srebrny, dżuma, samochód, t{-}shirt, spódnica} \\
            }
        }
    }
    & \\ \\

    \theutterance \stepcounter{utterance}  
    & & & \multicolumn{4}{p{0.6\linewidth}}{
        \cellcolor[rgb]{0.9,0.9,0.9}{
            \makecell[{{p{\linewidth}}}]{
                \texttt{\tiny{[GM$\langle$P2]}}
                \texttt{ZGADNIĘCIE: jurta} \\
            }
        }
    }
     \\ \\

    \theutterance \stepcounter{utterance}  
    & & & \multicolumn{2}{p{0.3\linewidth}}{
        \cellcolor[rgb]{0.9,0.9,0.9}{
            \makecell[{{p{\linewidth}}}]{
                \texttt{\tiny{[GM$|$GM]}}
                ['jurta']
            }
        }
    }
    & & \\ \\

    \theutterance \stepcounter{utterance}  
    & & & \multicolumn{2}{p{0.3\linewidth}}{
        \cellcolor[rgb]{0.9,0.9,0.9}{
            \makecell[{{p{\linewidth}}}]{
                \texttt{\tiny{[GM$|$GM]}}
                {'word': 'jurta', 'assignment': 'team'}
            }
        }
    }
    & & \\ \\

    \theutterance \stepcounter{utterance}  
    & & & \multicolumn{2}{p{0.3\linewidth}}{
        \cellcolor[rgb]{0.9,0.9,0.9}{
            \makecell[{{p{\linewidth}}}]{
                \texttt{\tiny{[GM$|$GM]}}
                {'word': 'jurta', 'assignment': 'team'}
            }
        }
    }
    & & \\ \\

    \theutterance \stepcounter{utterance}  
    & & & \multicolumn{2}{p{0.3\linewidth}}{
        \cellcolor[rgb]{0.9,0.9,0.9}{
            \makecell[{{p{\linewidth}}}]{
                \texttt{\tiny{[GM$|$GM]}}
                {'word': 'boks', 'assignment': 'opponent'}
            }
        }
    }
    & & \\ \\

    \theutterance \stepcounter{utterance}  
    & & & \multicolumn{2}{p{0.3\linewidth}}{
        \cellcolor[rgb]{0.9,0.9,0.9}{
            \makecell[{{p{\linewidth}}}]{
                \texttt{\tiny{[GM$|$GM]}}
                {'hidden': {'team': ['srebrny', 'irys', 'spódnica', 'chata', 't-shirt', 'narcyz'], 'innocent': ['sufit', 'gruźlica', 'bez', 'jaskinia', 'pokój', 'jeansy', 'dżuma'], 'opponent': ['ściana', 'biały', 'samochód', 'dach', 'majtki'], 'assassin': ['porwanie']}, 'revealed': {'team': {'team': ['tenis', 'sprzeniewierzenie', 'jurta'], 'innocent': [], 'opponent': [], 'assassin': []}, 'opponent': {'team': [], 'innocent': [], 'opponent': ['przekroczenie prędkości', 'petunia', 'boks'], 'assassin': []}}}
            }
        }
    }
    & & \\ \\

    \theutterance \stepcounter{utterance}  
    & & \multicolumn{4}{p{0.6\linewidth}}{
        \cellcolor[rgb]{0.9,0.9,0.9}{
            \makecell[{{p{\linewidth}}}]{
                \texttt{\tiny{[P1$\langle$GM]}}
                \texttt{Słowo jurta zostało odgadnięte prawidłowo. Tura twojego kolegi z drużyny się tam skończyła.} \\
\texttt{Teraz podaj kolejną wskazówkę odnoszącą się do niektórych z pozostałych słów twojego zespołu oraz listę powiązanych słów docelowych. Pamiętaj, aby rozpocząć wskazówkę od 'WSKAZÓWKA: ', wstawić nową linię i rozpocząć listę słów docelowych oddzielonych przecinkami od 'CELE: '. Uwaga: niektóre słowa zostały usunięte z list w porównaniu do poprzednich żądań.} \\
\\ 
\texttt{Twoje pozostałe słowa zespołu to: srebrny, irys, spódnica, chata, t{-}shirt, narcyz.} \\
\texttt{Pozostałe słowa twojego przeciwnika to: ściana, biały, samochód, boks, dach, majtki.} \\
\texttt{Pozostałe słowa rozpraszające to: sufit, gruźlica, bez, jaskinia, pokój, jeansy, dżuma.} \\
\texttt{Pozostałe słowa zabójcy to: porwanie.} \\
            }
        }
    }
    & \\ \\

    \theutterance \stepcounter{utterance}  
    & \multicolumn{4}{p{0.6\linewidth}}{
        \cellcolor[rgb]{0.9,0.9,0.9}{
            \makecell[{{p{\linewidth}}}]{
                \texttt{\tiny{[P1$\rangle$GM]}}
                \texttt{WSKAZÓWKA: mitologia} \\
\texttt{CELE: irys, narcyz} \\
            }
        }
    }
    & & \\ \\

    \theutterance \stepcounter{utterance}  
    & & & \multicolumn{2}{p{0.3\linewidth}}{
        \cellcolor[rgb]{0.9,0.9,0.9}{
            \makecell[{{p{\linewidth}}}]{
                \texttt{\tiny{[GM$|$GM]}}
                \texttt{mitologia} \\
            }
        }
    }
    & & \\ \\

    \theutterance \stepcounter{utterance}  
    & & & \multicolumn{2}{p{0.3\linewidth}}{
        \cellcolor[rgb]{0.9,0.9,0.9}{
            \makecell[{{p{\linewidth}}}]{
                \texttt{\tiny{[GM$|$GM]}}
                ['irys', 'narcyz']
            }
        }
    }
    & & \\ \\

    \theutterance \stepcounter{utterance}  
    & & & \multicolumn{2}{p{0.3\linewidth}}{
        \cellcolor[rgb]{0.9,0.9,0.9}{
            \makecell[{{p{\linewidth}}}]{
                \texttt{\tiny{[GM$|$GM]}}
                {'word': 'irys', 'assignment': 'team'}
            }
        }
    }
    & & \\ \\

    \theutterance \stepcounter{utterance}  
    & & & \multicolumn{2}{p{0.3\linewidth}}{
        \cellcolor[rgb]{0.9,0.9,0.9}{
            \makecell[{{p{\linewidth}}}]{
                \texttt{\tiny{[GM$|$GM]}}
                {'word': 'narcyz', 'assignment': 'team'}
            }
        }
    }
    & & \\ \\

    \theutterance \stepcounter{utterance}  
    & & \multicolumn{4}{p{0.6\linewidth}}{
        \cellcolor[rgb]{0.9,0.9,0.9}{
            \makecell[{{p{\linewidth}}}]{
                \texttt{\tiny{[GM$\rangle$P2]}}
                \texttt{Słowo jurta zostało odgadnięte prawidłowo. Twoja tura się tam skończyła.} \\
\texttt{Teraz podaj kolejną listę oddzieloną przecinkami, zawierającą co najmniej 1 i maksymalnie 2 słów z poniższej listy słów, które najlepiej odnoszą się lub są najbardziej związane ze słowem 'mitologia'. Pamiętaj, aby rozpocząć odpowiedź od ZGADNIĘCIE: . Uwaga: niektóre słowa zostały usunięte z listy w porównaniu do poprzednich próśb.} \\
\\ 
\texttt{jaskinia, majtki, dach, ściana, porwanie, pokój, chata, jeansy, sufit, biały, irys, gruźlica, narcyz, bez, srebrny, dżuma, samochód, t{-}shirt, spódnica} \\
            }
        }
    }
    & \\ \\

    \theutterance \stepcounter{utterance}  
    & & & \multicolumn{4}{p{0.6\linewidth}}{
        \cellcolor[rgb]{0.9,0.9,0.9}{
            \makecell[{{p{\linewidth}}}]{
                \texttt{\tiny{[GM$\langle$P2]}}
                \texttt{ZGADNIĘCIE: narcyz, irys} \\
            }
        }
    }
     \\ \\

    \theutterance \stepcounter{utterance}  
    & & & \multicolumn{2}{p{0.3\linewidth}}{
        \cellcolor[rgb]{0.9,0.9,0.9}{
            \makecell[{{p{\linewidth}}}]{
                \texttt{\tiny{[GM$|$GM]}}
                ['narcyz', 'irys']
            }
        }
    }
    & & \\ \\

    \theutterance \stepcounter{utterance}  
    & & & \multicolumn{2}{p{0.3\linewidth}}{
        \cellcolor[rgb]{0.9,0.9,0.9}{
            \makecell[{{p{\linewidth}}}]{
                \texttt{\tiny{[GM$|$GM]}}
                {'word': 'narcyz', 'assignment': 'team'}
            }
        }
    }
    & & \\ \\

    \theutterance \stepcounter{utterance}  
    & & & \multicolumn{2}{p{0.3\linewidth}}{
        \cellcolor[rgb]{0.9,0.9,0.9}{
            \makecell[{{p{\linewidth}}}]{
                \texttt{\tiny{[GM$|$GM]}}
                {'word': 'narcyz', 'assignment': 'team'}
            }
        }
    }
    & & \\ \\

    \theutterance \stepcounter{utterance}  
    & & & \multicolumn{2}{p{0.3\linewidth}}{
        \cellcolor[rgb]{0.9,0.9,0.9}{
            \makecell[{{p{\linewidth}}}]{
                \texttt{\tiny{[GM$|$GM]}}
                {'word': 'irys', 'assignment': 'team'}
            }
        }
    }
    & & \\ \\

    \theutterance \stepcounter{utterance}  
    & & & \multicolumn{2}{p{0.3\linewidth}}{
        \cellcolor[rgb]{0.9,0.9,0.9}{
            \makecell[{{p{\linewidth}}}]{
                \texttt{\tiny{[GM$|$GM]}}
                {'word': 'irys', 'assignment': 'team'}
            }
        }
    }
    & & \\ \\

    \theutterance \stepcounter{utterance}  
    & & & \multicolumn{2}{p{0.3\linewidth}}{
        \cellcolor[rgb]{0.9,0.9,0.9}{
            \makecell[{{p{\linewidth}}}]{
                \texttt{\tiny{[GM$|$GM]}}
                {'word': 'dach', 'assignment': 'opponent'}
            }
        }
    }
    & & \\ \\

    \theutterance \stepcounter{utterance}  
    & & & \multicolumn{2}{p{0.3\linewidth}}{
        \cellcolor[rgb]{0.9,0.9,0.9}{
            \makecell[{{p{\linewidth}}}]{
                \texttt{\tiny{[GM$|$GM]}}
                {'hidden': {'team': ['srebrny', 'spódnica', 'chata', 't-shirt'], 'innocent': ['sufit', 'gruźlica', 'bez', 'jaskinia', 'pokój', 'jeansy', 'dżuma'], 'opponent': ['ściana', 'biały', 'samochód', 'majtki'], 'assassin': ['porwanie']}, 'revealed': {'team': {'team': ['tenis', 'sprzeniewierzenie', 'jurta', 'narcyz', 'irys'], 'innocent': [], 'opponent': [], 'assassin': []}, 'opponent': {'team': [], 'innocent': [], 'opponent': ['przekroczenie prędkości', 'petunia', 'boks', 'dach'], 'assassin': []}}}
            }
        }
    }
    & & \\ \\

    \theutterance \stepcounter{utterance}  
    & & \multicolumn{4}{p{0.6\linewidth}}{
        \cellcolor[rgb]{0.9,0.9,0.9}{
            \makecell[{{p{\linewidth}}}]{
                \texttt{\tiny{[P1$\langle$GM]}}
                \texttt{Słowa narcyz, irys zostały odgadnięte prawidłowo. Tura twojego kolegi z drużyny się tam skończyła.} \\
\texttt{Teraz podaj kolejną wskazówkę odnoszącą się do niektórych z pozostałych słów twojego zespołu oraz listę powiązanych słów docelowych. Pamiętaj, aby rozpocząć wskazówkę od 'WSKAZÓWKA: ', wstawić nową linię i rozpocząć listę słów docelowych oddzielonych przecinkami od 'CELE: '. Uwaga: niektóre słowa zostały usunięte z list w porównaniu do poprzednich żądań.} \\
\\ 
\texttt{Twoje pozostałe słowa zespołu to: srebrny, spódnica, chata, t{-}shirt.} \\
\texttt{Pozostałe słowa twojego przeciwnika to: ściana, biały, samochód, dach, majtki.} \\
\texttt{Pozostałe słowa rozpraszające to: sufit, gruźlica, bez, jaskinia, pokój, jeansy, dżuma.} \\
\texttt{Pozostałe słowa zabójcy to: porwanie.} \\
            }
        }
    }
    & \\ \\

    \theutterance \stepcounter{utterance}  
    & \multicolumn{4}{p{0.6\linewidth}}{
        \cellcolor[rgb]{0.9,0.9,0.9}{
            \makecell[{{p{\linewidth}}}]{
                \texttt{\tiny{[P1$\rangle$GM]}}
                \texttt{WSKAZÓWKA: falbana} \\
\texttt{CELE: spódnica} \\
            }
        }
    }
    & & \\ \\

    \theutterance \stepcounter{utterance}  
    & & & \multicolumn{2}{p{0.3\linewidth}}{
        \cellcolor[rgb]{0.9,0.9,0.9}{
            \makecell[{{p{\linewidth}}}]{
                \texttt{\tiny{[GM$|$GM]}}
                \texttt{falbana} \\
            }
        }
    }
    & & \\ \\

    \theutterance \stepcounter{utterance}  
    & & & \multicolumn{2}{p{0.3\linewidth}}{
        \cellcolor[rgb]{0.9,0.9,0.9}{
            \makecell[{{p{\linewidth}}}]{
                \texttt{\tiny{[GM$|$GM]}}
                ['spódnica']
            }
        }
    }
    & & \\ \\

    \theutterance \stepcounter{utterance}  
    & & & \multicolumn{2}{p{0.3\linewidth}}{
        \cellcolor[rgb]{0.9,0.9,0.9}{
            \makecell[{{p{\linewidth}}}]{
                \texttt{\tiny{[GM$|$GM]}}
                {'word': 'spódnica', 'assignment': 'team'}
            }
        }
    }
    & & \\ \\

    \theutterance \stepcounter{utterance}  
    & & \multicolumn{4}{p{0.6\linewidth}}{
        \cellcolor[rgb]{0.9,0.9,0.9}{
            \makecell[{{p{\linewidth}}}]{
                \texttt{\tiny{[GM$\rangle$P2]}}
                \texttt{Słowa narcyz, irys zostały odgadnięte prawidłowo. Twoja tura się tam skończyła.} \\
\texttt{Teraz podaj kolejną listę oddzieloną przecinkami, zawierającą co najmniej 1 i maksymalnie 1 słów z poniższej listy słów, które najlepiej odnoszą się lub są najbardziej związane ze słowem 'falbana'. Pamiętaj, aby rozpocząć odpowiedź od ZGADNIĘCIE: . Uwaga: niektóre słowa zostały usunięte z listy w porównaniu do poprzednich próśb.} \\
\\ 
\texttt{jaskinia, majtki, ściana, porwanie, pokój, chata, jeansy, sufit, biały, gruźlica, bez, srebrny, dżuma, samochód, t{-}shirt, spódnica} \\
            }
        }
    }
    & \\ \\

    \theutterance \stepcounter{utterance}  
    & & & \multicolumn{4}{p{0.6\linewidth}}{
        \cellcolor[rgb]{0.9,0.9,0.9}{
            \makecell[{{p{\linewidth}}}]{
                \texttt{\tiny{[GM$\langle$P2]}}
                \texttt{ZGADNIĘCIE: spódnica} \\
            }
        }
    }
     \\ \\

    \theutterance \stepcounter{utterance}  
    & & & \multicolumn{2}{p{0.3\linewidth}}{
        \cellcolor[rgb]{0.9,0.9,0.9}{
            \makecell[{{p{\linewidth}}}]{
                \texttt{\tiny{[GM$|$GM]}}
                ['spódnica']
            }
        }
    }
    & & \\ \\

    \theutterance \stepcounter{utterance}  
    & & & \multicolumn{2}{p{0.3\linewidth}}{
        \cellcolor[rgb]{0.9,0.9,0.9}{
            \makecell[{{p{\linewidth}}}]{
                \texttt{\tiny{[GM$|$GM]}}
                {'word': 'spódnica', 'assignment': 'team'}
            }
        }
    }
    & & \\ \\

    \theutterance \stepcounter{utterance}  
    & & & \multicolumn{2}{p{0.3\linewidth}}{
        \cellcolor[rgb]{0.9,0.9,0.9}{
            \makecell[{{p{\linewidth}}}]{
                \texttt{\tiny{[GM$|$GM]}}
                {'word': 'spódnica', 'assignment': 'team'}
            }
        }
    }
    & & \\ \\

    \theutterance \stepcounter{utterance}  
    & & & \multicolumn{2}{p{0.3\linewidth}}{
        \cellcolor[rgb]{0.9,0.9,0.9}{
            \makecell[{{p{\linewidth}}}]{
                \texttt{\tiny{[GM$|$GM]}}
                {'word': 'samochód', 'assignment': 'opponent'}
            }
        }
    }
    & & \\ \\

    \theutterance \stepcounter{utterance}  
    & & & \multicolumn{2}{p{0.3\linewidth}}{
        \cellcolor[rgb]{0.9,0.9,0.9}{
            \makecell[{{p{\linewidth}}}]{
                \texttt{\tiny{[GM$|$GM]}}
                {'hidden': {'team': ['srebrny', 'chata', 't-shirt'], 'innocent': ['sufit', 'gruźlica', 'bez', 'jaskinia', 'pokój', 'jeansy', 'dżuma'], 'opponent': ['ściana', 'biały', 'majtki'], 'assassin': ['porwanie']}, 'revealed': {'team': {'team': ['tenis', 'sprzeniewierzenie', 'jurta', 'narcyz', 'irys', 'spódnica'], 'innocent': [], 'opponent': [], 'assassin': []}, 'opponent': {'team': [], 'innocent': [], 'opponent': ['przekroczenie prędkości', 'petunia', 'boks', 'dach', 'samochód'], 'assassin': []}}}
            }
        }
    }
    & & \\ \\

    \theutterance \stepcounter{utterance}  
    & & \multicolumn{4}{p{0.6\linewidth}}{
        \cellcolor[rgb]{0.9,0.9,0.9}{
            \makecell[{{p{\linewidth}}}]{
                \texttt{\tiny{[P1$\langle$GM]}}
                \texttt{Słowo spódnica zostało odgadnięte prawidłowo. Tura twojego kolegi z drużyny się tam skończyła.} \\
\texttt{Teraz podaj kolejną wskazówkę odnoszącą się do niektórych z pozostałych słów twojego zespołu oraz listę powiązanych słów docelowych. Pamiętaj, aby rozpocząć wskazówkę od 'WSKAZÓWKA: ', wstawić nową linię i rozpocząć listę słów docelowych oddzielonych przecinkami od 'CELE: '. Uwaga: niektóre słowa zostały usunięte z list w porównaniu do poprzednich żądań.} \\
\\ 
\texttt{Twoje pozostałe słowa zespołu to: srebrny, chata, t{-}shirt.} \\
\texttt{Pozostałe słowa twojego przeciwnika to: ściana, biały, samochód, majtki.} \\
\texttt{Pozostałe słowa rozpraszające to: sufit, gruźlica, bez, jaskinia, pokój, jeansy, dżuma.} \\
\texttt{Pozostałe słowa zabójcy to: porwanie.} \\
            }
        }
    }
    & \\ \\

    \theutterance \stepcounter{utterance}  
    & \multicolumn{4}{p{0.6\linewidth}}{
        \cellcolor[rgb]{0.9,0.9,0.9}{
            \makecell[{{p{\linewidth}}}]{
                \texttt{\tiny{[P1$\rangle$GM]}}
                \texttt{WSKAZÓWKA: kruszec} \\
\texttt{CELE: srebrny} \\
            }
        }
    }
    & & \\ \\

    \theutterance \stepcounter{utterance}  
    & & & \multicolumn{2}{p{0.3\linewidth}}{
        \cellcolor[rgb]{0.9,0.9,0.9}{
            \makecell[{{p{\linewidth}}}]{
                \texttt{\tiny{[GM$|$GM]}}
                \texttt{kruszec} \\
            }
        }
    }
    & & \\ \\

    \theutterance \stepcounter{utterance}  
    & & & \multicolumn{2}{p{0.3\linewidth}}{
        \cellcolor[rgb]{0.9,0.9,0.9}{
            \makecell[{{p{\linewidth}}}]{
                \texttt{\tiny{[GM$|$GM]}}
                ['srebrny']
            }
        }
    }
    & & \\ \\

    \theutterance \stepcounter{utterance}  
    & & & \multicolumn{2}{p{0.3\linewidth}}{
        \cellcolor[rgb]{0.9,0.9,0.9}{
            \makecell[{{p{\linewidth}}}]{
                \texttt{\tiny{[GM$|$GM]}}
                {'word': 'srebrny', 'assignment': 'team'}
            }
        }
    }
    & & \\ \\

    \theutterance \stepcounter{utterance}  
    & & \multicolumn{4}{p{0.6\linewidth}}{
        \cellcolor[rgb]{0.9,0.9,0.9}{
            \makecell[{{p{\linewidth}}}]{
                \texttt{\tiny{[GM$\rangle$P2]}}
                \texttt{Słowo spódnica zostało odgadnięte prawidłowo. Twoja tura się tam skończyła.} \\
\texttt{Teraz podaj kolejną listę oddzieloną przecinkami, zawierającą co najmniej 1 i maksymalnie 1 słów z poniższej listy słów, które najlepiej odnoszą się lub są najbardziej związane ze słowem 'kruszec'. Pamiętaj, aby rozpocząć odpowiedź od ZGADNIĘCIE: . Uwaga: niektóre słowa zostały usunięte z listy w porównaniu do poprzednich próśb.} \\
\\ 
\texttt{jaskinia, majtki, ściana, porwanie, pokój, chata, jeansy, sufit, biały, gruźlica, bez, srebrny, dżuma, t{-}shirt} \\
            }
        }
    }
    & \\ \\

    \theutterance \stepcounter{utterance}  
    & & & \multicolumn{4}{p{0.6\linewidth}}{
        \cellcolor[rgb]{0.9,0.9,0.9}{
            \makecell[{{p{\linewidth}}}]{
                \texttt{\tiny{[GM$\langle$P2]}}
                \texttt{ZGADNIĘCIE: srebrny} \\
            }
        }
    }
     \\ \\

    \theutterance \stepcounter{utterance}  
    & & & \multicolumn{2}{p{0.3\linewidth}}{
        \cellcolor[rgb]{0.9,0.9,0.9}{
            \makecell[{{p{\linewidth}}}]{
                \texttt{\tiny{[GM$|$GM]}}
                ['srebrny']
            }
        }
    }
    & & \\ \\

    \theutterance \stepcounter{utterance}  
    & & & \multicolumn{2}{p{0.3\linewidth}}{
        \cellcolor[rgb]{0.9,0.9,0.9}{
            \makecell[{{p{\linewidth}}}]{
                \texttt{\tiny{[GM$|$GM]}}
                {'word': 'srebrny', 'assignment': 'team'}
            }
        }
    }
    & & \\ \\

    \theutterance \stepcounter{utterance}  
    & & & \multicolumn{2}{p{0.3\linewidth}}{
        \cellcolor[rgb]{0.9,0.9,0.9}{
            \makecell[{{p{\linewidth}}}]{
                \texttt{\tiny{[GM$|$GM]}}
                {'word': 'srebrny', 'assignment': 'team'}
            }
        }
    }
    & & \\ \\

    \theutterance \stepcounter{utterance}  
    & & & \multicolumn{2}{p{0.3\linewidth}}{
        \cellcolor[rgb]{0.9,0.9,0.9}{
            \makecell[{{p{\linewidth}}}]{
                \texttt{\tiny{[GM$|$GM]}}
                {'word': 'majtki', 'assignment': 'opponent'}
            }
        }
    }
    & & \\ \\

    \theutterance \stepcounter{utterance}  
    & & & \multicolumn{2}{p{0.3\linewidth}}{
        \cellcolor[rgb]{0.9,0.9,0.9}{
            \makecell[{{p{\linewidth}}}]{
                \texttt{\tiny{[GM$|$GM]}}
                {'hidden': {'team': ['chata', 't-shirt'], 'innocent': ['sufit', 'gruźlica', 'bez', 'jaskinia', 'pokój', 'jeansy', 'dżuma'], 'opponent': ['ściana', 'biały'], 'assassin': ['porwanie']}, 'revealed': {'team': {'team': ['tenis', 'sprzeniewierzenie', 'jurta', 'narcyz', 'irys', 'spódnica', 'srebrny'], 'innocent': [], 'opponent': [], 'assassin': []}, 'opponent': {'team': [], 'innocent': [], 'opponent': ['przekroczenie prędkości', 'petunia', 'boks', 'dach', 'samochód', 'majtki'], 'assassin': []}}}
            }
        }
    }
    & & \\ \\

    \theutterance \stepcounter{utterance}  
    & & \multicolumn{4}{p{0.6\linewidth}}{
        \cellcolor[rgb]{0.9,0.9,0.9}{
            \makecell[{{p{\linewidth}}}]{
                \texttt{\tiny{[P1$\langle$GM]}}
                \texttt{Słowo srebrny zostało odgadnięte prawidłowo. Tura twojego kolegi z drużyny się tam skończyła.} \\
\texttt{Teraz podaj kolejną wskazówkę odnoszącą się do niektórych z pozostałych słów twojego zespołu oraz listę powiązanych słów docelowych. Pamiętaj, aby rozpocząć wskazówkę od 'WSKAZÓWKA: ', wstawić nową linię i rozpocząć listę słów docelowych oddzielonych przecinkami od 'CELE: '. Uwaga: niektóre słowa zostały usunięte z list w porównaniu do poprzednich żądań.} \\
\\ 
\texttt{Twoje pozostałe słowa zespołu to: chata, t{-}shirt.} \\
\texttt{Pozostałe słowa twojego przeciwnika to: ściana, biały, majtki.} \\
\texttt{Pozostałe słowa rozpraszające to: sufit, gruźlica, bez, jaskinia, pokój, jeansy, dżuma.} \\
\texttt{Pozostałe słowa zabójcy to: porwanie.} \\
            }
        }
    }
    & \\ \\

    \theutterance \stepcounter{utterance}  
    & \multicolumn{4}{p{0.6\linewidth}}{
        \cellcolor[rgb]{0.9,0.9,0.9}{
            \makecell[{{p{\linewidth}}}]{
                \texttt{\tiny{[P1$\rangle$GM]}}
                \texttt{WSKAZÓWKA: podkoszulek} \\
\texttt{CELE: t{-}shirt} \\
            }
        }
    }
    & & \\ \\

    \theutterance \stepcounter{utterance}  
    & & & \multicolumn{2}{p{0.3\linewidth}}{
        \cellcolor[rgb]{0.9,0.9,0.9}{
            \makecell[{{p{\linewidth}}}]{
                \texttt{\tiny{[GM$|$GM]}}
                \texttt{podkoszulek} \\
            }
        }
    }
    & & \\ \\

    \theutterance \stepcounter{utterance}  
    & & & \multicolumn{2}{p{0.3\linewidth}}{
        \cellcolor[rgb]{0.9,0.9,0.9}{
            \makecell[{{p{\linewidth}}}]{
                \texttt{\tiny{[GM$|$GM]}}
                ['t-shirt']
            }
        }
    }
    & & \\ \\

    \theutterance \stepcounter{utterance}  
    & & & \multicolumn{2}{p{0.3\linewidth}}{
        \cellcolor[rgb]{0.9,0.9,0.9}{
            \makecell[{{p{\linewidth}}}]{
                \texttt{\tiny{[GM$|$GM]}}
                {'word': 't-shirt', 'assignment': 'team'}
            }
        }
    }
    & & \\ \\

    \theutterance \stepcounter{utterance}  
    & & \multicolumn{4}{p{0.6\linewidth}}{
        \cellcolor[rgb]{0.9,0.9,0.9}{
            \makecell[{{p{\linewidth}}}]{
                \texttt{\tiny{[GM$\rangle$P2]}}
                \texttt{Słowo srebrny zostało odgadnięte prawidłowo. Twoja tura się tam skończyła.} \\
\texttt{Teraz podaj kolejną listę oddzieloną przecinkami, zawierającą co najmniej 1 i maksymalnie 1 słów z poniższej listy słów, które najlepiej odnoszą się lub są najbardziej związane ze słowem 'podkoszulek'. Pamiętaj, aby rozpocząć odpowiedź od ZGADNIĘCIE: . Uwaga: niektóre słowa zostały usunięte z listy w porównaniu do poprzednich próśb.} \\
\\ 
\texttt{jaskinia, ściana, porwanie, pokój, chata, jeansy, sufit, biały, gruźlica, bez, dżuma, t{-}shirt} \\
            }
        }
    }
    & \\ \\

    \theutterance \stepcounter{utterance}  
    & & & \multicolumn{4}{p{0.6\linewidth}}{
        \cellcolor[rgb]{0.9,0.9,0.9}{
            \makecell[{{p{\linewidth}}}]{
                \texttt{\tiny{[GM$\langle$P2]}}
                \texttt{ZGADNIĘCIE: t{-}shirt} \\
            }
        }
    }
     \\ \\

    \theutterance \stepcounter{utterance}  
    & & & \multicolumn{2}{p{0.3\linewidth}}{
        \cellcolor[rgb]{0.9,0.9,0.9}{
            \makecell[{{p{\linewidth}}}]{
                \texttt{\tiny{[GM$|$GM]}}
                ['t-shirt']
            }
        }
    }
    & & \\ \\

    \theutterance \stepcounter{utterance}  
    & & & \multicolumn{2}{p{0.3\linewidth}}{
        \cellcolor[rgb]{0.9,0.9,0.9}{
            \makecell[{{p{\linewidth}}}]{
                \texttt{\tiny{[GM$|$GM]}}
                {'word': 't-shirt', 'assignment': 'team'}
            }
        }
    }
    & & \\ \\

    \theutterance \stepcounter{utterance}  
    & & & \multicolumn{2}{p{0.3\linewidth}}{
        \cellcolor[rgb]{0.9,0.9,0.9}{
            \makecell[{{p{\linewidth}}}]{
                \texttt{\tiny{[GM$|$GM]}}
                {'word': 't-shirt', 'assignment': 'team'}
            }
        }
    }
    & & \\ \\

    \theutterance \stepcounter{utterance}  
    & & & \multicolumn{2}{p{0.3\linewidth}}{
        \cellcolor[rgb]{0.9,0.9,0.9}{
            \makecell[{{p{\linewidth}}}]{
                \texttt{\tiny{[GM$|$GM]}}
                {'word': 'ściana', 'assignment': 'opponent'}
            }
        }
    }
    & & \\ \\

    \theutterance \stepcounter{utterance}  
    & & & \multicolumn{2}{p{0.3\linewidth}}{
        \cellcolor[rgb]{0.9,0.9,0.9}{
            \makecell[{{p{\linewidth}}}]{
                \texttt{\tiny{[GM$|$GM]}}
                {'hidden': {'team': ['chata'], 'innocent': ['sufit', 'gruźlica', 'bez', 'jaskinia', 'pokój', 'jeansy', 'dżuma'], 'opponent': ['biały'], 'assassin': ['porwanie']}, 'revealed': {'team': {'team': ['tenis', 'sprzeniewierzenie', 'jurta', 'narcyz', 'irys', 'spódnica', 'srebrny', 't-shirt'], 'innocent': [], 'opponent': [], 'assassin': []}, 'opponent': {'team': [], 'innocent': [], 'opponent': ['przekroczenie prędkości', 'petunia', 'boks', 'dach', 'samochód', 'majtki', 'ściana'], 'assassin': []}}}
            }
        }
    }
    & & \\ \\

    \theutterance \stepcounter{utterance}  
    & & \multicolumn{4}{p{0.6\linewidth}}{
        \cellcolor[rgb]{0.9,0.9,0.9}{
            \makecell[{{p{\linewidth}}}]{
                \texttt{\tiny{[P1$\langle$GM]}}
                \texttt{Słowo t{-}shirt zostało odgadnięte prawidłowo. Tura twojego kolegi z drużyny się tam skończyła.} \\
\texttt{Teraz podaj kolejną wskazówkę odnoszącą się do niektórych z pozostałych słów twojego zespołu oraz listę powiązanych słów docelowych. Pamiętaj, aby rozpocząć wskazówkę od 'WSKAZÓWKA: ', wstawić nową linię i rozpocząć listę słów docelowych oddzielonych przecinkami od 'CELE: '. Uwaga: niektóre słowa zostały usunięte z list w porównaniu do poprzednich żądań.} \\
\\ 
\texttt{Twoje pozostałe słowa zespołu to: chata.} \\
\texttt{Pozostałe słowa twojego przeciwnika to: ściana, biały.} \\
\texttt{Pozostałe słowa rozpraszające to: sufit, gruźlica, bez, jaskinia, pokój, jeansy, dżuma.} \\
\texttt{Pozostałe słowa zabójcy to: porwanie.} \\
            }
        }
    }
    & \\ \\

    \theutterance \stepcounter{utterance}  
    & \multicolumn{4}{p{0.6\linewidth}}{
        \cellcolor[rgb]{0.9,0.9,0.9}{
            \makecell[{{p{\linewidth}}}]{
                \texttt{\tiny{[P1$\rangle$GM]}}
                \texttt{WSKAZÓWKA: chałupa} \\
\texttt{CELE: chata} \\
            }
        }
    }
    & & \\ \\

    \theutterance \stepcounter{utterance}  
    & & & \multicolumn{2}{p{0.3\linewidth}}{
        \cellcolor[rgb]{0.9,0.9,0.9}{
            \makecell[{{p{\linewidth}}}]{
                \texttt{\tiny{[GM$|$GM]}}
                \texttt{chałupa} \\
            }
        }
    }
    & & \\ \\

    \theutterance \stepcounter{utterance}  
    & & & \multicolumn{2}{p{0.3\linewidth}}{
        \cellcolor[rgb]{0.9,0.9,0.9}{
            \makecell[{{p{\linewidth}}}]{
                \texttt{\tiny{[GM$|$GM]}}
                ['chata']
            }
        }
    }
    & & \\ \\

    \theutterance \stepcounter{utterance}  
    & & & \multicolumn{2}{p{0.3\linewidth}}{
        \cellcolor[rgb]{0.9,0.9,0.9}{
            \makecell[{{p{\linewidth}}}]{
                \texttt{\tiny{[GM$|$GM]}}
                {'word': 'chata', 'assignment': 'team'}
            }
        }
    }
    & & \\ \\

    \theutterance \stepcounter{utterance}  
    & & \multicolumn{4}{p{0.6\linewidth}}{
        \cellcolor[rgb]{0.9,0.9,0.9}{
            \makecell[{{p{\linewidth}}}]{
                \texttt{\tiny{[GM$\rangle$P2]}}
                \texttt{Słowo t{-}shirt zostało odgadnięte prawidłowo. Twoja tura się tam skończyła.} \\
\texttt{Teraz podaj kolejną listę oddzieloną przecinkami, zawierającą co najmniej 1 i maksymalnie 1 słów z poniższej listy słów, które najlepiej odnoszą się lub są najbardziej związane ze słowem 'chałupa'. Pamiętaj, aby rozpocząć odpowiedź od ZGADNIĘCIE: . Uwaga: niektóre słowa zostały usunięte z listy w porównaniu do poprzednich próśb.} \\
\\ 
\texttt{jaskinia, porwanie, pokój, chata, jeansy, sufit, biały, gruźlica, bez, dżuma} \\
            }
        }
    }
    & \\ \\

    \theutterance \stepcounter{utterance}  
    & & & \multicolumn{4}{p{0.6\linewidth}}{
        \cellcolor[rgb]{0.9,0.9,0.9}{
            \makecell[{{p{\linewidth}}}]{
                \texttt{\tiny{[GM$\langle$P2]}}
                \texttt{ZGADNIĘCIE: chata} \\
            }
        }
    }
     \\ \\

    \theutterance \stepcounter{utterance}  
    & & & \multicolumn{2}{p{0.3\linewidth}}{
        \cellcolor[rgb]{0.9,0.9,0.9}{
            \makecell[{{p{\linewidth}}}]{
                \texttt{\tiny{[GM$|$GM]}}
                ['chata']
            }
        }
    }
    & & \\ \\

    \theutterance \stepcounter{utterance}  
    & & & \multicolumn{2}{p{0.3\linewidth}}{
        \cellcolor[rgb]{0.9,0.9,0.9}{
            \makecell[{{p{\linewidth}}}]{
                \texttt{\tiny{[GM$|$GM]}}
                {'word': 'chata', 'assignment': 'team'}
            }
        }
    }
    & & \\ \\

    \theutterance \stepcounter{utterance}  
    & & & \multicolumn{2}{p{0.3\linewidth}}{
        \cellcolor[rgb]{0.9,0.9,0.9}{
            \makecell[{{p{\linewidth}}}]{
                \texttt{\tiny{[GM$|$GM]}}
                {'word': 'chata', 'assignment': 'team'}
            }
        }
    }
    & & \\ \\

    \theutterance \stepcounter{utterance}  
    & & & \multicolumn{2}{p{0.3\linewidth}}{
        \cellcolor[rgb]{0.9,0.9,0.9}{
            \makecell[{{p{\linewidth}}}]{
                \texttt{\tiny{[GM$|$GM]}}
                \texttt{team has won} \\
            }
        }
    }
    & & \\ \\

\end{supertabular}
}

\end{document}
